\documentclass[11pt]{article}
\usepackage[margin=0.8in]{geometry}
\usepackage{booktabs}
\usepackage{longtable}
\usepackage{array}
\usepackage{amsmath}
\usepackage{enumitem}
\usepackage{xcolor}
\usepackage[hidelinks]{hyperref}
\usepackage{microtype}

\setlength{\parindent}{0pt}
\setlength{\parskip}{5pt}
\setlist{nosep,leftmargin=1.4em}
\newcommand{\code}[1]{\texttt{#1}}

\title{\textbf{CS 121 Search Engine}\\Milestone 3 Report}
\author{Kevin Guo}
\date{June 3, 2026}

\begin{document}
\maketitle

\section{System Overview}

This project implements a search engine for the full ICS developer corpus without
using a text-indexing library. It consists of two separate Python programs:

\begin{itemize}
  \item \code{indexer.py} parses the crawled JSON documents and builds a
  disk-based inverted index.
  \item \code{searcher.py} accepts console queries, retrieves postings from disk,
  calculates TF-IDF-based scores, and returns ranked URLs.
\end{itemize}

The Milestone 3 retrieval component uses ranked retrieval rather than Boolean
AND retrieval. A document containing any available query term can receive a
score. Documents matching multiple query terms accumulate more TF-IDF score and
receive an additional coverage benefit, but a document is not excluded merely
because it is missing one query term.

\section{Indexer Design}

\subsection{Document Processing}

Each corpus document is read from JSON using its \code{url} and \code{content}
fields. URL fragments are removed before the URL is stored. HTML is parsed with
BeautifulSoup, and script, style, and noscript content is excluded. Plain-text
documents and unusually large HTML documents use a lighter extraction path to
avoid pathological parsing time while still indexing their textual content.

Tokens are all alphanumeric sequences matched by
\code{[A-Za-z0-9]+}. Tokens are lowercased and Porter stemmed. Stop words and
numeric tokens are retained to follow the no-stopping and all-alphanumeric-token
requirements.

\subsection{Important Words}

Text appearing in \code{title}, \code{h1}, \code{h2}, \code{h3}, \code{b}, and
\code{strong} tags is counted separately. Each posting stores:

\begin{center}
\code{(document ID, term frequency, important-term frequency)}
\end{center}

The searcher applies an important-word multiplier during scoring so occurrences
in these tags contribute more strongly than ordinary body text.

\subsection{Partial Indexes and Merge}

The in-memory inverted index is flushed to a sorted partial index every 5,000
documents. On the full corpus this produces more than the required three
offloads. After indexing, a k-way heap merge combines the partial indexes into
one sorted \code{index.txt} file.

During the merge, the indexer records the byte offset of each token in
\code{index\_seek.json}. It also writes \code{doc\_map.json}, which maps numeric
document IDs to URLs, and \code{analytics.json}, which summarizes the index.

\section{Search and Ranking}

\subsection{Disk-Based Retrieval}

The searcher loads the compact seek table and document-to-URL map, but it does
not load the full inverted index into memory. For each stemmed query term, it
uses the term's byte offset to seek directly to the corresponding postings line
in \code{index.txt}.

\subsection{TF-IDF Ranked Retrieval}

Candidate documents are the union of the retrieved postings lists. For each
query term $t$ and document $d$, the base contribution is:

\[
  \operatorname{score}(t,d) =
  \left(1 + \log(tf_{t,d} + 5 \cdot important_{t,d})\right)
  \left(\log\left(\frac{N+1}{df_t+1}\right)+1\right)
\]

when the effective term frequency is positive. Here, $N$ is the number of
indexed documents and $df_t$ is the document frequency of the term. Scores are
summed across all query terms found in the document.

This is not Boolean AND retrieval: a document matching one query term can be
ranked alongside a document matching all terms. Matching more query terms
normally increases the accumulated TF-IDF score. A small coverage multiplier
also favors documents that match a larger fraction of the query.

\subsection{Additional Ranking Components}

\begin{itemize}
  \item Query terms found in a URL receive a bonus.
  \item Raw text/data URLs and dataset paths receive mild penalties because
  repeated terms in large data files can otherwise dominate results.
  \item Equivalent URL forms such as \code{/}, \code{/index}, and
  \code{/index.html} are collapsed in the final result list.
  \item URL-based reranking is limited to a preliminary high-scoring window to
  control query time.
\end{itemize}

\section{Index Analytics}

\begin{center}
\begin{tabular}{lr}
\toprule
Metric & Value \\
\midrule
Indexed documents & 55,393 \\
Unique tokens & 388,597 \\
Inverted-index size & 93,216.91 KB (approximately 91.0 MB) \\
\bottomrule
\end{tabular}
\end{center}

\section{Evaluation}

The final searcher was evaluated on 20 queries after the seek table and document
map were loaded. Every query completed within the developer-track 300 ms target.

\begin{center}
\begin{longtable}{>{\raggedright\arraybackslash}p{0.48\linewidth}r}
\toprule
Query & Time (ms) \\
\midrule
\endfirsthead
\toprule
Query & Time (ms) \\
\midrule
\endhead
\code{cristina lopes} & 7.7 \\
\code{machine learning} & 87.8 \\
\code{ACM} & 92.7 \\
\code{student affairs} & 105.3 \\
\code{artificial intelligence} & 67.2 \\
\code{data mining} & 88.0 \\
\code{donald bren school} & 72.6 \\
\code{secure systems} & 107.8 \\
\code{software engineering masters} & 131.1 \\
\code{informatics graduate program} & 93.1 \\
\code{master of software engineering} & 185.1 \\
\code{database systems} & 115.1 \\
\code{computer vision} & 77.2 \\
\code{mondego lab} & 78.6 \\
\code{undergraduate admissions} & 80.6 \\
\code{ics faculty} & 102.0 \\
\code{course catalogue} & 102.1 \\
\code{networked systems} & 78.1 \\
\code{game design} & 87.7 \\
\code{uci calendar} & 128.0 \\
\midrule
Average & 94.4 \\
Maximum & 185.1 \\
Queries under 300 ms & 20 of 20 \\
\bottomrule
\end{longtable}
\end{center}

The move from Boolean AND to union-based TF-IDF retrieval improved recall:
pages may now be returned even when they contain only part of a multi-term
query. The tested results were generally relevant, and URL-aware ranking helped
surface program, calendar, faculty, and course pages. Some broad queries still
rank publication pages or dataset pages above the most authoritative page,
showing that TF-IDF alone does not fully model page quality.

\section{Operational Requirements}

\begin{itemize}
  \item The index is stored in filesystem files rather than a database.
  \item The indexer offloads its in-memory index to disk repeatedly and merges
  the partial indexes at the end.
  \item The searcher performs direct disk seeks and does not load all postings
  into memory.
  \item The search interface is an interactive console prompt and identifies
  results by URL.
  \item The measured response time remained below 300 ms for all 20 evaluation
  queries after startup.
\end{itemize}

\section{Limitations and Future Improvements}

\begin{itemize}
  \item Ranking quality could be improved with document-length normalization
  such as cosine normalization or BM25.
  \item Content-based exact and near-duplicate detection would reduce duplicate
  pages beyond the current URL normalization.
  \item Positional indexes and phrase matching would improve multi-word queries.
  \item PageRank, anchor text, or domain-quality signals could help authoritative
  pages outrank repeated-term publication and dataset pages.
\end{itemize}

\section{Running the System}

\begin{verbatim}
python3 indexer.py developer.zip -o output
python3 searcher.py --index-dir output
python3 searcher.py --index-dir output --query "machine learning" --top 5
\end{verbatim}

\end{document}
